\begin{recipe}{Pasta bolognese}
\kicker[Hoofdgerecht,Vlees,Italiaans,Doordeweeks]{Hoofdgerecht · vlees · Italiaans · doordeweeks}
\index[register]{Pasta bolognese}
\index[register]{Gehakt}
\index[register]{Soffrito}
\index[register]{Tomaat}
\index[register]{Tomatenpuree}

\meta{90 minuten}

\lettrine{D}{e} kinderen noemen dit pasta rode saus, hoewel het zeker niet de
enige rode saus in huis is. De saus draait om een fijn gesneden soffrito van
bleekselderij, winterpeen en ui in gelijke hoeveelheden, die naast het gehakt
voor body zorgt. Ik maak hem het liefst in een grote pan, in porties in te
vriezen voor een drukke avond.

\blockrule{Ingrediënten}
\begin{ingredients}
  \ing{200 g}{half-om-half gehakt}\index[register]{Gehakt}
  \ing{350 g}{tagliatelle (of andere brede lintpasta)}\index[register]{Pasta}
  \ing{3 stengels}{bleekselderij}
  \ing{1}{winterpeen}
  \ing{1 grote}{ui}
  \ing{1 blik (400\,g)}{gepelde tomaten (een goed merk maakt hier merkbaar verschil)}\index[register]{Tomaat}
  \ing{70 g}{tomatenpuree}
  \ing{1}{runderbouillonblokje}
  \ing{50 g}{Parmezaanse kaas, geraspt}\index[register]{Parmezaanse kaas}
  \ingb{scheut witte wijn, voor afblussen}
  \ingb{olijfolie}
  \ingb{versgemalen peper en zout}
\end{ingredients}

\blockrule{Bereiding}
\tip{Verdeel de saus in porties van 750 gram en laat op een drukke dag
vanuit de vriezer opwarmen in een hapjespan, terwijl de pasta kookt. Ragù
hoort op brede lintpasta of in lasagne; spaghetti houdt de saus minder goed
vast.}
\begin{steps}
  \item Schil de winterpeen en snijd samen met de bleekselderij en de ui
        heel fijn, in gelijke hoeveelheden: dit heet soffrito, of mirepoix
        in de Franse keuken. Verhit een eetlepel olijfolie op laag vuur in
        een hapjespan en voeg de groente toe.
  \item Laat de soffrito op laag vuur bakken tot de ui helemaal glazig is,
        na 10 tot 15 minuten: de groente hoeft niet bruin te worden. Haal de
        soffrito uit de pan.\index[register]{Soffrito}
  \item Zet het vuur flink hoger. Bak het gehakt rul in dezelfde pan met het
        bouillonblokje tot het mooi bruin gebakken is.
  \item Voeg de tomatenpuree toe en bak circa 1 minuut mee onder goed
        roeren: dit ontzuurt de puree.
  \item Blus af met een scheut witte wijn tot het vocht grotendeels is
        verdampt. Voeg de tomaten toe met een half blik water.
  \item Voeg de soffrito weer toe en roer goed door.
  \item Doe het deksel op de pan, zet het vuur laag en laat de saus minstens
        een uur zachtjes sudderen. Te nat? Laat hem nog even zonder deksel
        doorgaan. Te droog? Voeg een scheutje water toe.
  \item Voor de vriezer: laat de saus afkoelen tot kamertemperatuur voor je
        hem in porties verdeelt over ziplockzakjes.
  \item Kook de pasta al dente. Meng de saus door de pasta, eventueel met een
        scheutje kookvocht, en roer er ruim Parmezaanse kaas door.
\end{steps}
\heroimagefade{images/pasta-bolognese}
\end{recipe}
